\documentclass[12pt,a4paper]{article}

% =============================================================================
% PAQUETES
% =============================================================================
\usepackage[utf8]{inputenc}
\usepackage[T1]{fontenc}
\usepackage[spanish]{babel}
\usepackage{geometry}
\usepackage{graphicx}
\usepackage{xcolor}
\usepackage{booktabs}
\usepackage{longtable}
\usepackage{array}
\usepackage{multirow}
\usepackage{float}
\usepackage{caption}
\usepackage{subcaption}
\usepackage{amsmath}
\usepackage{amssymb}
\usepackage{enumitem}
\usepackage{fancyhdr}
\usepackage{titlesec}
\usepackage{setspace}
\usepackage{tikz}
\usepackage{pgfplots}
\pgfplotsset{compat=1.18}

% =============================================================================
% CONFIGURACION DE PAGINA
% =============================================================================
\geometry{
    left=2.5cm,
    right=2.5cm,
    top=2.5cm,
    bottom=2.5cm
}

% =============================================================================
% COLORES
% =============================================================================
\definecolor{azuloscuro}{RGB}{0, 51, 102}
\definecolor{azulclaro}{RGB}{70, 130, 180}
\definecolor{azulclarito}{RGB}{163,186,240}
\definecolor{grisclaro}{RGB}{245, 245, 245}
\definecolor{rojoacento}{RGB}{178, 34, 34}

% =============================================================================
% CONFIGURACION DE ENCABEZADOS
% =============================================================================
\pagestyle{fancy}
\fancyhf{}
\fancyhead[L]{\textcolor{azuloscuro}{\small Análisis de Datos en Ingeniería I}}
\fancyhead[R]{\textcolor{azuloscuro}{\small Proyecto Final - RLM}}
\fancyfoot[C]{\thepage}
\renewcommand{\headrulewidth}{0.5pt}
\renewcommand{\footrulewidth}{0.3pt}

% =============================================================================
% CONFIGURACION DE TITULOS
% =============================================================================
\titleformat{\section}
    {\normalfont\Large\bfseries\color{azuloscuro}}
    {\thesection.}{0.5em}{}
    
\titleformat{\subsection}
    {\normalfont\large\bfseries\color{azulclaro}}
    {\thesubsection}{0.5em}{}

\titleformat{\subsubsection}
    {\normalfont\large\bfseries\color{azulclarito}}
    {\thesubsection}{0.5em}{}

% =============================================================================
% ESPACIADO
% =============================================================================
\setstretch{1.5}

% =============================================================================
% INICIO DEL DOCUMENTO
% =============================================================================
\begin{document}

% =============================================================================
% PORTADA
% =============================================================================
\begin{titlepage}

    
    \vspace*{0.7cm}
    
    \begin{center}
        {\Huge\bfseries\textcolor{azuloscuro}{PROYECTO FINAL}}\\[0.5cm]
        {\LARGE\bfseries\textcolor{azulclaro}{Análisis de Datos en Ingeniería I}}\\[1cm]
        
        \rule{0.8\textwidth}{1.5pt}\\[0.8cm]
        
        {\Large\textbf{Regresión Lineal Múltiple:}}\\[0.3cm]
        {\Large\textbf{Predicción de Precios de Laptops}}\\[1.5cm]
        
        \rule{0.6\textwidth}{0.5pt}\\[1cm]
        
        {\large\textbf{Integrantes:}}\\[0.5cm]
        \begin{tabular}{rl}
            Dubal Aguilar Torres & (200207378) \\
            Alejandro Chaves Ramos & (200116116) \\
            Juan Caceres Figueroa & (200079919) \\
            Miguel Carrizosa Morales & (200208123) \\
        \end{tabular}
        
        \vfill
        
        {\large\textbf{Docente:}}\\[0.3cm]
        PhD. Luis Angel Anillo Arrieta\\[1cm]
        
        {\large Ingeniería de Sistemas}\\[0.3cm]
        {\large Universidad del Norte}\\[0.3cm]
        {\large Barranquilla, Colombia}\\[0.5cm]
        {\large 25 de mayo de 2026}
        
    \end{center}
\end{titlepage}

\newpage
\tableofcontents
\newpage


\section{Introducción}

En la actualidad, el mercado de computadores portátiles ha experimentado un crecimiento exponencial. Estos dispositivos representan una inversión importante para los consumidores, con precios que varían ampliamente según sus especificaciones técnicas o marca. Comprender los factores que determinan el precio de una laptop es importante tanto para consumidores como para fabricantes y distribuidores.

Desde la perspectiva de la ingeniería de sistemas, el análisis de datos permite identificar patrones y relaciones entre las características técnicas de un equipo y su precio en el mercado. El uso de estas técnicas estadísticas, como lo puede ser la regresión lineal múltiple permite cuantificar el impacto de variables específicas como pueden ser variables como la memoria RAM, la frecuencia del procesador, el tamaño de pantalla, el tipo de laptop o el sistema operativo, sobre el precio final del producto, para poder tomar la mejor decisión respecto a su compra.

Para este proyecto se utiliza el Laptop Price Prediction Dataset, un dataset que recopila información sobre las especificaciones técnicas y precios de 1303 laptops de diversas características. A partir de estos datos, se busca estudiar y evaluar qué variables influyen en el precio de mercado y en qué medida lo hacen. Esto permitirá comprender mejor el comportamiento del mercado de hardware y ofrecer una visión más clara de los elementos que determinan el valor de un equipo, ayudando tanto a consumidores como a empresas.

\subsection{Objetivo General}

Evaluar y predecir los factores técnicos y de categoría que influyen en el precio de las laptops, utilizando técnicas de regresión lineal múltiple.

\subsection{Objetivos Específicos}

\begin{enumerate}[leftmargin=2cm]
    \item Organizar el conjunto de datos para identificar las variables numéricas y categóricas relevantes asociadas al precio de laptops.
    
    \item Aplicar un modelo de regresión lineal múltiple para evaluar la relación entre las variables seleccionadas y el precio del equipo.
    
    \item Interpretar los coeficientes del modelo para determinar cuáles variables tienen el mayor impacto en el precio y qué implicaciones tienen para la toma de decisiones en la compra de estos equipos.
\end{enumerate}



\section{Plan de Análisis}

El análisis se desarrolló siguiendo la metodología de regresión lineal múltiple. Primero se seleccionaron las variables cuantitativas y cualitativas del dataset relacionadas con el precio de laptops. Después se exploraron sus relaciones, se evaluó la correlación, se generaron gráficos y finalmente se estimó el modelo de regresión.

Para validar los resultados, se verificaron los supuestos de normalidad, homocedasticidad, independencia y ausencia de multicolinealidad. Todo el proceso siguió los pasos metodológicos que se vieron y desarrollaron durante el curso.

En este proyecto se utiliza el modelo con un propósito predictivo, es decir, el objetivo es predecir el precio de una laptop a partir de las variables seleccionadas.

\subsection{Selección de Variables}

Se eligieron variables que, según la lógica del mercado de hardware, pueden influir sobre la variable dependiente:

\begin{itemize}[leftmargin=2cm]
    \item \textbf{Y (dependiente):} Precio de la laptop (euros)
    
    \item \textbf{Variables cuantitativas:}
    \begin{itemize}
        \item X1: Memoria RAM (GB)
        \item X2: Frecuencia del CPU (GHz)
        \item X3: Tamaño de la pantalla (pulgadas)
    \end{itemize}
    
    \item \textbf{Variables cualitativas:}
    \begin{itemize}
        \item X4: Tipo de laptop (Ultrabook, Notebook, Gaming, Workstation, etc.)
        \item X5: Sistema operativo (Windows, macOS, Linux, etc.)
    \end{itemize}
\end{itemize}

Estas variables fueron seleccionadas porque presentan correlaciones aceptables con la variable dependiente. La correlación de Pearson entre el precio y la RAM es de 0.74 fuerte, entre el precio y la frecuencia del CPU es de 0.43 moderada, y entre el precio y el tamaño de pantalla es de 0.07 débil pero presente.


\section{Análisis de Datos}
En esta sección el grupo realizará el análisis de los datos tanto cualitativos como cuantitativos y buscará entender, además de explicar las relaciones entre las variables dependientes e independientes, a través de la estadística, y pruebas para estimar parámetros relevantes.

\subsection{Descripción del Dataset}

El dataset utilizado fue el de Laptop Price Prediction Dataset, disponible en la plataforma Kaggle. Contiene información sobre 1303 laptops de diferentes categorías, incluyendo especificaciones técnicas y precios en euros.

Para este análisis, se utilizaron todos los registros disponibles (1,303 laptops) con las siguientes características:

\begin{table}[H]
\centering
\caption{Resumen del Dataset}
\begin{tabular}{lc}
\toprule
\textbf{Característica} & \textbf{Valor} \\
\midrule
Total de laptops & 1,303 \\
Variables numéricas & 3 (RAM, CPU GHz, Pulgadas) \\
Variables cualitativas & 2 (Tipo, SO) \\
Variable dependiente & Precio en euros \\
\bottomrule
\end{tabular}
\end{table}

\subsection{Correlación de Pearson}

La matriz de correlación de Pearson nos permite identificar las relaciones lineales entre las variables numéricas del conjunto de datos. A continuación se presenta la matriz de correlación para las variables seleccionadas:

\begin{table}[H]
\centering
\caption{Matriz de Correlación de Pearson}
\begin{tabular}{lcccc}
\toprule
 & \textbf{Precio} & \textbf{RAM (GB)} & \textbf{CPU (GHz)} & \textbf{Pulgadas} \\
\midrule
Precio & 1.0000 & 0.7430 & 0.4303 & 0.0682 \\
RAM (GB) & 0.7430 & 1.0000 & 0.3680 & 0.2380 \\
CPU (GHz) & 0.4303 & 0.3680 & 1.0000 & 0.3079 \\
Pulgadas & 0.0682 & 0.2380 & 0.3079 & 1.0000 \\
\bottomrule
\end{tabular}
\end{table}

Interpretación:
\begin{itemize}
    \item La correlación entre Precio y RAM es \textbf{0.743}, lo que indica una relación fuerte y positiva. A medida que aumenta la memoria RAM, también se aumentará el precio de la laptop.
    
    \item La correlación entre Precio y CPU (GHz) es \textbf{0.430}, lo que indica una relación moderada y positiva. Las laptops con procesadores de mayor frecuencia tienden a tener precios más altos.
    
    \item La correlación entre Precio y Pulgadas es \textbf{0.068}, lo que indica una relación muy débil. El tamaño de pantalla por sí solo no determina fuertemente el precio de la laptop.
    
\end{itemize}

\subsubsection{Gráficas Relación entre variables cuantitativas y precio}
Para analizar los distintos factores que afectan el precio de las laptops, se realizaron representaciones gráficas que permiten demostrar de manera visual el comportamiento de las distintas variables cuantitativas.
\vspace{5mm}
\begin{figure}[H]
    \centering
    \includegraphics[width=0.5\textwidth]{img/RAM.png}
    \caption{Gráfico RAM vs Precio.}
    \label{fig:RAM}
\end{figure}
\vspace{2mm}
En esta figura se puede observar la alta relación entre la cantidad de RAM y el precio de las laptops, se puede observar que tienen una relación creciente, además de que tienen una correlación fuerte, la razón de que los valores se vean en columnas y no sigan completamente la línea, se debe al estándar del mercado en donde las RAM se construyen en potencias de 2, o sea 2 de RAM 4 de RAM 8 de RAM y siguiendo ese formato. 
\vspace{5mm}
\begin{figure}[H]
    \centering
    \includegraphics[width=0.5\textwidth]{img/frecuenciaCPU.png}
    \caption{Gráfico CPU vs Precio.}
    \label{fig:CPU}
\end{figure}
\vspace{2mm}
En esta gráfica se puede ver la relación existente entre la frecuencia de la CPU y el precio de la laptop, se observa que tiene una relación creciente, y que cuenta con una correlación moderada, por lo que se puede entender que la frecuencia de la CPU tiene menos incidencia en el precio de la laptop que factores como lo pueden ser la RAM.

\vspace{5mm}
\begin{figure}[H]
    \centering
    \includegraphics[width=0.5\textwidth]{img/tamaPantalla.png}
    \caption{Gráfico Pulgadas vs Precio.}
    \label{fig:pulgadas}
\end{figure}
\vspace{2mm}
En esta figura observamos la relación que se da entre las pulgadas de la pantalla y los precios de los equipos, en este caso podemos observar que esta es la variable que menos afecta el precio de las laptops esto concuerda con el modelo que se planteó, además, al igual que pasaba con la RAM se ubican en columnas, que se ven explicadas por los estándares de la industria.
\subsection{Relaciones entre Variables Cualitativas y Precio}

Para analizar la relación entre las variables cualitativas y el precio, se utilizaron diagramas de caja y bigotes que nos permiten visualizar la distribución del precio según cada categoría.


\begin{figure}[H]
    \centering
    \includegraphics[width=0.5\textwidth]{img/prec_tipo.png}
    \caption{Diagrama de caja que muestra la distribución del precio según el tipo de laptop (Ultrabook, Notebook, Gaming, Workstation, etc.).}
    \label{fig:tipo}
\end{figure}
\vspace{2mm}
En la figura 4, se evidencia la relación existente entre los tipos de laptop y el precio de las mismas. Se puede observar que cada una de las categorías presentan distribuciones distintas, esto mostrándonos que el tipo de laptop genera diferencias en el precio.

Analizando con detalle la gráfica podemos observar que basándonos en el promedio el tipo de laptop que tiene un precio mayor son las del tipo WorkStation, y las de menor precio siendo las del tipo Netbook.

Por otro lado, se puede observar que se presentaron datos atípicos que presentan precios mayores que el promedio por tipo de laptop, estos aumentos de precios extremos no están asociados a un tipo de laptop específico, sino, a los diversos factores que afectan al precio de una laptop como se mencionó con anterioridad. Estos resultados son consistentes con el modelo de regresión que se presentará a continuación afirmando que el tipo de laptop, es determinante en el precio de la laptop.
\vspace{5mm}

% ESPACIO PARA FIGURA 3: BOXPLOT OpSys vs PRECIO
\begin{figure}[H]
    \centering
    \includegraphics[width=0.5\textwidth]{img/prec_ope.png}
    \caption{Diagrama de caja que muestra la distribución del precio según el sistema operativo (Windows, macOS, Linux, etc.).}
    \label{fig:opsys}
\end{figure}
\vspace{2mm}
En la figura 5, se evidencia la relación que existe entre los distintos sistemas operativos y el precio de las laptops. En estas gráficas se puede observar que cada una de las categorías presentan distribuciones distintas, esto dándonos a entender que el sistema operativo con el que venga la laptop afecta de manera significativa el precio de la misma.

Analizando con detenimiento la gráfica se observa que basándonos en el promedio de precio de cada uno de los sistemas operativos el sistema operativo que cuenta con los precios más alto es el de macOS y las de menor precio siendo las de Android.

Además, se puede observar que se presentaron datos atípicos que presentan precios mayores que el promedio por sistema operativos, estos aumentos de precios extremos no están asociados a un sistema operativo concreto, sino, a los diversos factores que afectan al precio de una laptop como se mencionó con anterioridad. Estos resultados son consistentes con el modelo de regresión que se presentará a continuación afirmando que el sistema operativo de una laptop, puede influir en el precio de la laptop.


\section{Modelo de Regresión Lineal Múltiple}

\subsection{Especificación del Modelo}

El modelo de regresión lineal múltiple estimado es:

\begin{equation}
    Precio_i = \beta_0 + \beta_1 \cdot RAM_i + \beta_2 \cdot CPU_i + \beta_3 \cdot Pulgadas_i + \beta_4 \cdot Tipo_i + \beta_5 \cdot SO_i + \varepsilon_i
\end{equation}

Donde:
\begin{itemize}
    \item $Precio_i$: Precio en euros de la laptop $i$
    \item $RAM_i$: Memoria RAM en GB
    \item $CPU_i$: Frecuencia del procesador en GHz
    \item $Pulgadas_i$: Tamaño de pantalla en pulgadas
    \item $Tipo_i$: Variables dummy para el tipo de laptop
    \item $SO_i$: Variables dummy para el sistema operativo
    \item $\varepsilon_i$: Término de error
\end{itemize}

\subsection{Resultados del Modelo}

\begin{table}[H]
\centering
\caption{Resultados del Modelo RLM - Coeficientes Numéricos}
\begin{tabular}{lcccc}
\toprule
\textbf{Variable} & \textbf{Coeficiente} & \textbf{Error Est.} & \textbf{t} & \textbf{p-valor} \\
\midrule
(Intercepto) & 42.920 & 292.168 & 0.147 & 0.883 \\
RAM (GB) & 81.075 & 2.581 & 31.414 & \textless{}2e-16 *** \\
CPU (GHz) & 208.313 & 24.493 & 8.505 & \textless{}2e-16 *** \\
Pulgadas & -23.088 & 11.150 & -2.071 & 0.039 * \\
\bottomrule
\end{tabular}
\\[0.3cm]
\small Significancia: *** p\textless{}0.001, ** p\textless{}0.01, * p\textless{}0.05
\end{table}

\begin{table}[H]
\centering
\caption{Resultados del Modelo RLM - Variables Cualitativas}
\begin{tabular}{lcccc}
\toprule
\textbf{Variable} & \textbf{Coeficiente} & \textbf{Error Est.} & \textbf{t} & \textbf{p-valor} \\
\midrule
TypeNameGaming & -46.892 & 53.627 & -0.874 & 0.382 \\
TypeNameNetbook & -247.928 & 88.427 & -2.804 & 0.005 ** \\
TypeNameNotebook & -293.233 & 42.865 & -6.841 & \textless{}2e-16 *** \\
TypeNameUltrabook & 120.528 & 45.988 & 2.621 & 0.009 ** \\
TypeNameWorkstation & 650.140 & 85.032 & 7.646 & \textless{}2e-16 *** \\
\midrule
OpSysChrome OS & 263.411 & 282.896 & 0.931 & 0.352 \\
OpSysLinux & 231.796 & 279.931 & 0.828 & 0.408 \\
OpSysMac OS X & 401.450 & 306.282 & 1.311 & 0.190 \\
OpSysmacOS & 617.621 & 295.898 & 2.087 & 0.037 * \\
OpSysNo OS & 179.496 & 279.594 & 0.642 & 0.521 \\
OpSysWindows 10 & 420.911 & 274.979 & 1.531 & 0.126 \\
OpSysWindows 10 S & 502.645 & 306.600 & 1.639 & 0.101 \\
OpSysWindows 7 & 827.967 & 280.866 & 2.948 & 0.003 ** \\
\bottomrule
\end{tabular}
\\[0.3cm]
\small Las categorías de referencia son: TypeName = ``2 in 1 Convertible'', OpSys = ``Android''
\end{table}

\begin{table}[H]
\centering
\caption{Bondad de Ajuste}
\begin{tabular}{lc}
\toprule
\textbf{Métrica} & \textbf{Valor} \\
\midrule
R-cuadrado & 0.706 \\
R-cuadrado Ajustado & 0.703 \\
F-statistic & 193.3 \\
p-valor (F) & \textless{}2.2e-16 \\
Error estándar residual & 381.2 \\
Grados de libertad (modelo) & 16 \\
Grados de libertad (residuos) & 1,286 \\
\bottomrule
\end{tabular}
\end{table}

El resultado del modelo es satisfactorio, pues explica aproximadamente el 70.6\% de la variabilidad del precio de una laptop, y los valores de nuestro estadístico F y P-Valor nos indican que nuestro modelo tiene una validez estadística del análisis con un error estándar de estimación de 381.2 EUR.

\subsection{Modelo}
El modelo planteado quedaría de la siguiente manera:
\begin{equation}
\begin{aligned}
\widehat{\text{Precio}}_i = & \ 42.920 + 81.075 \cdot \text{RAM}_i + 208.313 \cdot \text{CPU}_i - 23.088 \cdot \text{Pulgadas}_i \\
& - 46.892 \cdot \text{Gaming}_i - 247.928 \cdot \text{Netbook}_i - 293.233 \cdot \text{Notebook}_i \\
& + 120.528 \cdot \text{Ultrabook}_i + 650.140 \cdot \text{Workstation}_i \\
& + 263.411 \cdot \text{Chrome OS}_i + 231.796 \cdot \text{Linux}_i + 401.450 \cdot \text{Mac OS X}_i \\
& + 617.621 \cdot \text{macOS}_i + 179.496 \cdot \text{No OS}_i + 420.911 \cdot \text{Windows 10}_i \\
& + 502.645 \cdot \text{Windows 10 S}_i + 827.967 \cdot \text{Windows 7}_i
\end{aligned}
\end{equation}


\subsection{Interpretación de Coeficientes}

\begin{itemize}
    \item \textbf{RAM (GB):} Por cada aumento de 1 GB en la memoria RAM, el precio de la laptop aumenta en promedio 81.08 EUR, manteniendo constantes las demás variables.
    
    \item \textbf{CPU (GHz):} Por cada aumento de 1 GHz en la frecuencia del procesador, el precio aumenta en promedio 208.31 EUR. Esto indica que el procesador es uno de los componentes más determinantes del precio.
    
    \item \textbf{Pulgadas:} Por cada pulgada adicional en el tamaño de pantalla, el precio \textbf{disminuye} en promedio 23.09 EUR. Este resultado contraintuitivo puede explicarse por el hecho de que las laptops con pantallas más grandes tienden a ser notebooks básicas, mientras que las más caras (Ultrabooks, Workstations) suelen tener pantallas más compactas pero de mayor resolución.
    
    \item \textbf{Tipo de Laptop:} Las Workstation tienen un precio promedio 650.14 EUR mayor que las ``2 in 1 Convertible'' (categoría de referencia), mientras que los Notebook son 293.23 EUR más baratos. Los Netbook son 247.93 EUR más baratos.
    
    \item \textbf{Sistema Operativo:} Windows 7 tiene un precio promedio 827.97 EUR mayor que Android (categoría de referencia). Mac OS es 617.62 EUR más caro. Sin embargo, la mayoría de las categorías de SO no son estadísticamente significativas individualmente (solo macOS p=0.037 y Windows 7 p=0.003).
\end{itemize}
    
    
    A pesar de que algunos resultados puedan parecer contraintuitivos comparando con las gráficas, debemos tener en cuenta las características de las variables que estamos usando, un resultado que podría parecer confuso, es el de que en nuestro modelo un aumento en la CPU tenga un mayor peso que la RAM, pero esto se puede explicar con el hecho de que el costo de aumentar la frecuencia en una CPU es más costoso y complejo que el de aumentar la RAM, esto se puede ver si nos damos cuenta que la mayoría de CPUs modernas no superan los 3GHz. 
    
    El otro resultado que puede parecer confuso es el del coeficiente del sistema operativo de Windows 7 sea tan alto, factor que se explica por la dificultad de encontrar un pc con las mismas características que use Windows 7, y también que los precios de los laptops han disminuido en comparación a los momentos donde se usaba Windows 7.

    \subsection{Intervalos de confianza}
    
    \begin{table}[H]
    \centering
    \caption{Intervalos de Confianza (95\%) del Modelo RLM}
    \begin{tabular}{lcc}
    \toprule
    \textbf{Variable} & \textbf{2.5 \%} & \textbf{97.5 \%} \\
    \midrule
    (Intercept) & -530.259 & 616.098 \\
    Ram\_GB & 76.012 & 86.138 \\
    CPU\_GHz & 160.262 & 256.364 \\
    Inches & -44.963 & -1.214 \\
    \midrule
    TypeNameGaming & -152.098 & 58.314 \\
    TypeNameNetbook & -421.405 & -74.451 \\
    TypeNameNotebook & -377.326 & -209.141 \\
    TypeNameUltrabook & 30.309 & 210.747 \\
    TypeNameWorkstation & 483.324 & 816.957 \\
    \midrule
    OpSysChrome OS & -291.577 & 818.399 \\
    OpSysLinux & -317.376 & 780.968 \\
    OpSysMac OS X & -199.417 & 1002.317 \\
    OpSysmacOS & 37.126 & 1198.117 \\
    OpSysNo OS & -369.014 & 728.006 \\
    OpSysWindows 10 & -118.545 & 960.367 \\
    OpSysWindows 10 S & -98.846 & 1104.135 \\
    OpSysWindows 7 & 276.962 & 1378.973 \\
    \bottomrule
    \end{tabular}
    \\[0.3cm]
    \small Las categorías de referencia son: TypeName = ``2 in 1 Convertible'', OpSys = ``Android''
    \end{table}

    En el cuadro 6 se pueden observar los intervalos de confianza, los cuales nos permiten evaluar 
la precisión y significancia de los coeficientes que se estimaron en el modelo de regresión. Los 
intervalos de confianza de las variables cuantitativas RAM y CPU no incluyen al 0, lo que indica que estas variables tienen un efecto consistente sobre la variable dependiente. En el caso de Pulgadas, el intervalo no incluye el 0, lo que sugiere que su efecto es estadísticamente significativo al 95\% de confianza, aunque su magnitud sea pequeña. Respecto a las variables cualitativas, algunas categorías como Gaming, Chrome OS, Linux, Mac OS X, No OS, Windows 10 y Windows 10 S tienen intervalos que incluyen el 0, indicando que no son indicativos consistentes del precio de una laptop. 

    Algo importante a tener en cuenta es que en el intervalo de confianza del intercepto se incluye el 0, lo que significa que no podemos afirmar con certeza que el intercepto sea diferente de cero. Sin embargo, el intercepto no tiene interpretación práctica en este contexto, ya que una laptop con RAM=0, CPU=0GHz y pantalla de 0 pulgadas no existe en la realidad.


\section{Validación de los Supuestos del Modelo}

\subsection{1. Independencia de Residuos (Durbin-Watson)}

\begin{table}[H]
\centering
\begin{tabular}{lc}
\toprule
\textbf{Prueba} & \textbf{Resultado} \\
\midrule
H0 & Los residuos son independientes \\
H1 & Los residuos son dependientes \\
Estadístico DW & 2.0272 \\
p-valor & 0.633 \\
\bottomrule
\end{tabular}
\end{table}

Interpretación: Como el p-valor (0.633) es mayor que 0.05, no se rechaza H0. Los residuos son independientes. El valor del estadístico Durbin-Watson (2.0272) está cercano a 2, lo que confirma la ausencia de autocorrelación.

\subsection{2. Normalidad de Residuos (Shapiro-Wilk)}

\begin{table}[H]
\centering
\begin{tabular}{lc}
\toprule
\textbf{Prueba} & \textbf{Resultado} \\
\midrule
H0 & Los residuos se distribuyen normalmente \\
H1 & Los residuos NO se distribuyen normalmente \\
Estadístico W & 0.92379 \\
p-valor & \textless{}2.2e-16 \\
\bottomrule
\end{tabular}
\end{table}

Interpretación: Como el p-valor es muy pequeño (\textless{}2.2e-16), se rechaza H0. \textbf{Los residuos NO siguen una distribución normal.} Este resultado puede atribuirse a la asimetría de la variable precio, que presenta una cola larga hacia la derecha debido a laptops de muy alto precio.




\begin{figure}[H]
    \centering
    \includegraphics[width=0.5\textwidth]{img/QQ_plot_final.png}
    \caption{Gráfico Q-Q de los residuos para evaluar normalidad. Los puntos se desvían de la línea diagonal en los extremos, indicando que los residuos no siguen una distribución perfectamente normal.}
    \label{fig:qqplot}
\end{figure}



\subsection{3. Homocedasticidad (Breusch-Pagan)}

\begin{table}[H]
\centering
\begin{tabular}{lc}
\toprule
\textbf{Prueba} & \textbf{Resultado} \\
\midrule
H0 & Los residuos son homocedásticos \\
H1 & Los residuos son heterocedásticos \\
Estadístico BP & 201.21 \\
df & 16 \\
p-valor & \textless{}2.2e-16 \\
\bottomrule
\end{tabular}
\end{table}

Interpretación: Como el p-valor es muy pequeño (\textless{}2.2e-16), se rechaza H0. \textbf{Los residuos son heterocedásticos.} Esto indica que la varianza de los residuos no es constante a lo largo de los valores ajustados. Este resultado es común en datos donde el rango de la variable dependiente es amplio, precios desde 174 hasta 6,099 EUR.





\begin{figure}[H]
    \centering
    \includegraphics[width=0.5\textwidth]{img/resi_ajustados.png}   
    \caption{Gráfico de dispersión de residuos contra valores ajustados.}
    \label{fig:residuos}
\end{figure}

\subsection{4. Multicolinealidad (VIF)}

\begin{table}[H]
\centering
\caption{Factor de Inflación de Varianza (VIF)}
\begin{tabular}{lccl}
\toprule
\textbf{Variable} & \textbf{VIF} & \textbf{Estado} & \textbf{Interpretación} \\
\midrule
RAM (GB) & 1.54 & Aceptable & Sin multicolinealidad \\
CPU (GHz) & 1.38 & Aceptable & Sin multicolinealidad \\
Pulgadas & 2.27 & Aceptable & Sin multicolinealidad \\
TypeName (conjunto) & 3.79 & Aceptable & Sin multicolinealidad \\
OpSys (conjunto) & 1.60 & Aceptable & Sin multicolinealidad \\
\bottomrule
\end{tabular}
\end{table}

Interpretación: Los valores de VIF son todos menores a 5, lo que indica que \textbf{no hay multicolinealidad problemática} en el modelo. Las variables numéricas tienen VIF muy bajos (menores a 2.5), lo que confirma que las correlaciones entre variables independientes son aceptables.


\begin{figure}[H]
    \centering
    \includegraphics[width=0.5\textwidth]{img/diag_completo.png}   
    \caption{Conjunto de gráficos de diagnóstico: residuos vs fitted, Q-Q plot, scale-location y residuos vs leverage.}
    \label{fig:diagnostico}
\end{figure}
\begin{itemize}
    \item \textbf{Residuals vs Fitted.}

        La gráfica demuestra que los residuos se distribuyen cerca de cero, pero se observa mayor dispersión en los valores ajustados altos, lo que nos indica heterocedasticidad, es decir, varianza no constante en los errores. El modelo predice mejor en laptops de rango de precio medio, pero comete más errores en valores altos de precio.

    \item \textbf{Q-Q Plot.}
        
        La gráfica nos muestra que los residuos no siguen perfectamente una distribución normal, siendo en los valores extremos donde se encuentra principalmente esta desviación.

    \item \textbf{Scale-Location.}
        
        Esta gráfica nos demuestra que la varianza no es constante, esto respalda el resultado de la presencia de heterocedasticidad en nuestro modelo, a pesar de esto, el modelo sigue siendo 
        útil para la predicción.

    \item \textbf{Residuals vs Leverage.}

        En esta gráfica se puede observar que a pesar de los valores atípicos que se encuentran en el dataset, ningún valor supera los límites críticos de distancia de Cook, y como además la mayoría de los puntos están concentrados en la zona izquierda de la gráfica, no se afecta la estabilidad del modelo.
\end{itemize}


\section{Conclusiones específicas}

\begin{enumerate}[leftmargin=1.5cm]
    \item El modelo de regresión lineal múltiple desarrollado explica aproximadamente el 70.6\% de la variabilidad en el precio de las laptops (R² = 0.706), lo que indica un buen ajuste del modelo.
    
    \item La frecuencia del CPU es el factor cuantitativo más influyente en el precio de una laptop. Por cada aumento de 1 GHz en la frecuencia del procesador, el precio aumenta en promedio 208.31 EUR. La memoria RAM también es importante: por cada GB adicional, el precio aumenta 81.08 EUR.
    
    \item El tamaño de pantalla muestra un efecto negativo en el precio (-23.09 EUR por pulgada), lo que sugiere que las laptops más caras tienden a ser más compactas pero con especificaciones superiores en otros aspectos.
    
    \item Las categorías de tipo de laptop y sistema operativo muestran diferencias significativas en el precio, siendo las Workstation las de mayor precio promedio (650.14 EUR más que las ``2 in 1 Convertible'').
    
    \item El supuesto de independencia se cumple (Durbin-Watson = 2.0272, p = 0.633) y no hay multicolinealidad (VIF \textless{} 5).
    
\end{enumerate}

\subsection{Conclusión general}

En conclusión, se encontró que los factores que más afectan el precio de un ordenador portátil son la frecuencia de la CPU y la memoria RAM, ya que demuestran una fuerte relación lineal con el costo final del portátil, siendo los más importantes en la valoración comercial del hardware. Además, mediante el análisis de las variables cualitativas y el comportamiento de los datos atípicos presentados en las gráficas, se concluye que las desviaciones extremas corresponden a propias características del mercado como pueden ser equipos empresariales de alta gama o modelos de portátiles de muy alta gama, donde el costo se ve influenciado no solamente por las especificaciones del sistema sino también por elementos de diseño, exclusividad o licencias de sistemas operativos específicos como Windows 7. 

Por otro lado, el modelo de regresión lineal múltiple presentado demostró una buena capacidad predictiva al explicar aproximadamente el 70.6\% de la variabilidad de los precios. Si bien algunos supuestos estadísticos, tales como la normalidad de los residuos y la homocedasticidad, no se cumplieron perfectamente debido a la dispersión propia en los segmentos de gama alta, el modelo refleja de manera válida y adecuada las tendencias generales del conjunto de datos que se usó. Asimismo, este modelo puede utilizarse en el mundo corporativo y en la gestión de tecnologías de la información para la toma de decisiones informadas que permitan optimizar las estrategias de adquisición de hardware.


\subsection{Recomendaciones}

\begin{itemize}[leftmargin=1.5cm]
    \item Mejorar la elección del dataset, generará modelos más rígidos y funcionales.
    
    \item Considerar transformaciones de la variable dependiente o usar modelos robustos para manejar la heterocedasticidad y no normalidad de los residuos.
    
    \item Para este modelo en concreto se sugiere incluir más variables como tipo de almacenamiento (SSD vs HDD), marca de GPU, y resolución de pantalla para mejorar el poder explicativo del modelo, también realizar un mejor tratamiento de los datos.
\end{itemize}

\subsection{Aprendizajes del Curso}

Este proyecto permitió aplicar de manera integral los conceptos aprendidos en el curso de Análisis de Datos:

\begin{itemize}[leftmargin=1.5cm]
    \item La importancia de verificar los supuestos del modelo antes de interpretar los resultados.
    \item El manejo de datos reales que requieren limpieza y preparación previa.
    \item La interpretación práctica de coeficientes en contextos reales de ingeniería de sistemas.
    \item La aplicación de pruebas estadísticas para validar los supuestos de regresión.
    \item La necesidad de reconocer limitaciones del modelo y proponer mejoras.
\end{itemize}


\newpage
\section{Anexos}

\subsection{Anexo A: Base de Datos Utilizada}

La base de datos completa está disponible en el archivo \texttt{laptop\_price.csv}, que contiene 1,303 laptops con las variables utilizadas en el análisis.

Fuente original: Laptop Price Prediction Dataset, Kaggle.
URL: https://www.kaggle.com/datasets/muhammetvarl/laptop-price

\subsection{Anexo B: Resultados Completos del Modelo}

\begin{verbatim}
Call:
lm(formula = Price_euros ~ Ram_GB + CPU_GHz + Inches + TypeName + 
    OpSys, data = datos)

Residuals:
    Min      1Q  Median      3Q     Max 
-1895.37 -218.45  -46.23  159.84  2882.97 

Coefficients:
                    Estimate Std. Error t value Pr(>|t|)    
(Intercept)          42.920    292.168   0.147 0.883233    
Ram_GB               81.075      2.581  31.414  < 2e-16 ***
CPU_GHz             208.313     24.493   8.505  < 2e-16 ***
Inches              -23.088     11.150  -2.071 0.038590 *  
TypeNameGaming      -46.892     53.627  -0.874 0.382057    
TypeNameNetbook    -247.928     88.427  -2.804 0.005132 ** 
TypeNameNotebook   -293.233     42.865  -6.841 1.21e-11 ***
TypeNameUltrabook   120.528     45.988   2.621 0.008873 ** 
TypeNameWorkstation 650.140     85.032   7.646 4.04e-14 ***
OpSysChrome OS      263.411    282.896   0.931 0.351971    
OpSysLinux          231.796    279.931   0.828 0.407796    
OpSysMac OS X       401.450    306.282   1.311 0.190192    
OpSysmacOS          617.621    295.898   2.087 0.037063 *  
OpSysNo OS          179.496    279.594   0.642 0.520985    
OpSysWindows 10     420.911    274.979   1.531 0.126089    
OpSysWindows 10 S   502.645    306.600   1.639 0.101366    
OpSysWindows 7      827.967    280.866   2.948 0.003257 ** 
---
Signif. codes:  0 '***' 0.001 '**' 0.01 '*' 0.05 '.' 0.1 ' ' 1

Residual standard error: 381.2 on 1286 degrees of freedom
Multiple R-squared:  0.7063,    Adjusted R-squared:  0.7027 
F-statistic: 193.3 on 16 and 1286 DF,  p-value: < 2.2e-16
\end{verbatim}

\subsection{Anexo C: Resumen de las Pruebas de Supuestos}

\begin{verbatim}
Prueba de Shapiro-Wilk:
  W = 0.92379, p-value < 2.2e-16
  Resultado: NO NORMAL

Prueba de Jarque-Bera:
  X-squared = 2411.5, df = 2, p-value < 2.2e-16
  Resultado: NO NORMAL

Prueba de Durbin-Watson:
  DW = 2.0272, p-value = 0.6331
  Resultado: INDEPENDIENTES

Prueba de Breusch-Pagan:
  BP = 201.21, df = 16, p-value < 2.2e-16
  Resultado: HETEROCEDASTICO

VIF (Multicolinealidad):
  Ram_GB: 1.54 (OK)
  CPU_GHz: 1.38 (OK)
  Inches: 2.27 (OK)
  TypeName: 3.79 (OK)
  OpSys: 1.60 (OK)
\end{verbatim}

\end{document}
